% Generated draft: requires mathematical and source review.
\documentclass[11pt]{article}
\usepackage[margin=25mm]{geometry}
\usepackage{fontspec}
\setmainfont{Latin Modern Roman}
\usepackage{amsmath,amssymb,mathtools,unicode-math}
\setmathfont{Latin Modern Math}
\usepackage{enumitem,needspace,xurl,hyperref}
\setlength{\parindent}{0pt}
\setlength{\parskip}{5pt}
\setlength{\emergencystretch}{3em}
\newcommand{\N}{\mathbb{N}}
\newcommand{\Z}{\mathbb{Z}}
\newcommand{\Q}{\mathbb{Q}}
\newcommand{\R}{\mathbb{R}}
\newcommand{\C}{\mathbb{C}}
\newcommand{\F}{\mathbb{F}}
\newcommand{\A}{\mathbb{A}}
\newcommand{\PP}{\mathbb{P}}
\newcommand{\E}{\mathbb{E}}
\newcommand{\eps}{\varepsilon}
\newcommand{\dd}{\,\mathrm{d}}
\begin{document}
\section*{UnsolvedMath ID: 700250}
\textbf{Source problem code:} catalogue entry
\subsection{Definitions and mathematical statement}
For each positive integer \(n\), let \(M_n(R)\) denote the set of \(n\times n\) matrices over a ring \(R\), and put
\[
\begin{aligned}
GL_n(\R)&=\{A\in M_n(\R):\det A\ne0\},\\
GL_n(\Z)&=\{V\in M_n(\Z):\det V\in\{1,-1\}\},\\
D_n(\R)&=\{\operatorname{diag}(d_1,\ldots,d_n):d_i\in\R\setminus\{0\}\},\\
O_n(\R)&=\{Q\in M_n(\R):Q^{\mathsf T}Q=I_n\},\\
UT_n(\R)&=\{U=(u_{ij})\in M_n(\R):u_{ii}=1\text{ for every }i,
\ u_{ij}=0\text{ whenever }i>j\}.
\end{aligned}
\]
These definitions use the customary interpretation of \(D_n\) as the invertible diagonal group and \(UT_n\) as the upper unitriangular group. The supplied excerpt refers elsewhere for notation and does not explicitly define these symbols; confirmation against that reference remains necessary.

For subsets \(X_1,\ldots,X_4\subseteq GL_n(\R)\), their product means
\[
X_1X_2X_3X_4=\{x_1x_2x_3x_4:x_i\in X_i\text{ for }1\le i\le4\}.
\]
The problem proposed by Yu. I. Merzlyakov asks to determine exactly the positive integers \(n\) for which
\[
GL_n(\R)=D_n(\R)\,O_n(\R)\,UT_n(\R)\,GL_n(\Z).
\]
Equivalently, determine the positive integers \(n\) such that
\[
\forall A\in GL_n(\R)\quad
\exists D\in D_n(\R)\ \exists Q\in O_n(\R)\
\exists U\in UT_n(\R)\ \exists V\in GL_n(\Z)
\quad A=DQUV.
\]
The order of the four factors is part of the question. The source says ``natural numbers'' without specifying whether \(0\) is included; the formulation above covers positive dimensions, leaving any intended zero-dimensional convention unresolved.

The supplied record reports that the equality holds for \(1\le n\le3\), citing Kh. N. Narzullayev, and that it fails for every sufficiently large positive integer \(n\), citing N. S. Akhmedov. No numerical threshold for the latter assertion is supplied. It also reports that, for each fixed \(n\), this equality implies an affirmative solution of Minkowski's problem on the product of \(n\) linear forms, citing A. M. Macbeath, and describes that problem as remaining open for \(n\ge6\). These are supplied source qualifications, not independently verified mathematical or current literature assessments.
\subsection{Short English statement}
Determine exactly which positive dimensions \(n\) allow every invertible real \(n\times n\) matrix to be written, in order, as an invertible diagonal matrix, an orthogonal matrix, an upper unitriangular matrix, and an integer matrix of determinant \(\pm1\). The notation requires confirmation against the cited reference. The supplied record reports validity for \(n\le3\), failure in all sufficiently large dimensions, and a connection to Minkowski's problem on products of linear forms.
\subsection{Sources}
\begin{itemize}
\item Local UnsolvedMath registry, supplied complete source record: permanent ID \(700250\), exact source problem code \texttt{catalogue entry}, title \emph{Kourovka Notebook Problem 4.42}. The registry labels the problem ``open''; that label is supplied metadata and has not been independently verified. Record updated \texttt{2026-10-16T14:00:00Z}; this is a registry timestamp, not a supplied literature-verification date.
\item \emph{The Kourovka Notebook (Unsolved Problems in Group Theory)}, problem 4.42, originating in the 4th issue (1973), as attributed in the record to the 21st issue edition (2026), arXiv:\texttt{1401.0300v48}. Proposed by Yu. I. Merzlyakov. The supplied collection description names E. I. Khukhro and V. D. Mazurov as editors.
\item M. I. Kargapolov and Ju. I. Merzljakov, \emph{Fundamentals of the Theory of Groups}, Springer, New York, 1979. Supplied reference for notation; no page or section is specified.
\item A. M. Macbeath, \emph{Proc. Glasgow Math. Assoc.}, volume 5, no. 2 (1961), pp. 86--89. Supplied reference concerning Minkowski's problem on the product of \(n\) linear forms; no article title is supplied.
\item Kh. N. Narzullayev, \emph{Math. Notes}, volume 18 (1975), pp. 713--719. Supplied reference for the assertion that the equality holds for \(n\le3\); no article title is supplied.
\item N. S. Akhmedov, \emph{Zapiski Nauchn. Seminarov LOMI}, volume 67 (1977), pp. 86--107 (Russian). Supplied reference for failure in all sufficiently large dimensions; no article title or explicit dimension threshold is supplied.
\end{itemize}
\end{document}
